\section{句法结构的两种视角}

\subsection{短语结构语法（Phrase Structure Grammar）}

语言学中表示句子结构的第一种方式是\textbf{短语结构}（Phrase Structure），由\textbf{上下文无关文法}（Context-Free Grammar, CFG）来形式化描述。

\begin{knowledgebox}{上下文无关文法（CFG）基本概念}
上下文无关文法通过一组\textbf{重写规则}描述语言的结构。每条规则将一个非终结符展开为终结符或非终结符的序列。例如：
\begin{itemize}
    \item $\text{NP} \rightarrow \text{Det} \ (\text{Adj})^* \ \text{N}$（名词短语 = 限定词 + 可选的形容词 + 名词）
    \item $\text{PP} \rightarrow \text{P} \ \text{NP}$（介词短语 = 介词 + 名词短语）
    \item $\text{VP} \rightarrow \text{V} \ \text{PP}$（动词短语 = 动词 + 介词短语）
    \item $\text{NP} \rightarrow \text{NP} \ \text{PP}$（名词短语也可以递归包含介词短语）
\end{itemize}
\end{knowledgebox}

以短语 ``the cuddly cat by the door'' 为例，其短语结构树自底向上构建：

\begin{itemize}
    \item ``the cuddly cat'' 构成一个名词短语（NP）
    \item ``the door'' 构成另一个名词短语
    \item ``by the door'' 构成一个介词短语（PP）
    \item 整体 ``the cuddly cat by the door'' 是一个更大的名词短语
\end{itemize}

\begin{figure}[H]
\centering
\includegraphics[width=0.95\textwidth]{slides-images/slide-002.jpg}
\caption{短语结构语法示例：词类（parts of speech）与短语构成规则\protect\footnotemark}
\end{figure}
\footnotetext{Slides 第2页。}

\begin{figure}[H]
\centering
\includegraphics[width=0.95\textwidth]{slides-images/slide-003.jpg}
\caption{上下文无关文法规则示例：NP、PP、VP 的重写规则\protect\footnotemark}
\end{figure}
\footnotetext{Slides 第3页。}

\subsection{依存结构（Dependency Structure）}

本讲的重点是第二种表示方式——\textbf{依存结构}（Dependency Structure）。依存结构不关心短语的层次嵌套，而是直接刻画\textbf{词与词之间的修饰关系}。

\begin{importantbox}{依存结构的核心思想}
依存结构通过\textbf{有向弧}（directed arcs）连接句子中的词对，表示\textbf{中心词}（head）与\textbf{依存词}（dependent/modifier）之间的关系。例如在 ``look in the large crate in the kitchen by the door'' 中：
\begin{itemize}
    \item ``look'' 是整个句子的中心词（root）
    \item ``crate'' 是 ``in'' 的宾语，``in the large crate'' 修饰 ``look''
    \item ``large'' 和 ``the'' 修饰 ``crate''
    \item ``in the kitchen'' 和 ``by the door'' 也修饰 ``crate''
\end{itemize}
\end{importantbox}

\begin{figure}[H]
\centering
\includegraphics[width=0.95\textwidth]{slides-images/slide-004.jpg}
\caption{短语结构与依存结构的并列对比：同一句话的两种结构表示\protect\footnotemark}
\end{figure}
\footnotetext{Slides 第4页。}

\begin{figure}[H]
\centering
\includegraphics[width=0.95\textwidth]{slides-images/slide-005.jpg}
\caption{依存结构示意：从 ``look'' 出发，箭头指向各个依存词，形成树形结构\protect\footnotemark}
\end{figure}
\footnotetext{Slides 第5页。}

\subsection{本章小结}

句法结构有两种主要表示方式：短语结构（基于上下文无关文法的树形嵌套）和依存结构（基于词间修饰关系的有向图）。两者可以相互转化，但本课程重点使用依存结构，因为它更直接地表达词与词之间的语义关系，且在 NLP 实践中更为流行。

%% ========== 句法歧义 ==========

